%%%%%%%%%%%%%%%%%%%%%%%%%%%%%%%%%%%%%%%%%%%%%%%%%%%%%%%%%%%%%%%%%%%%%%%%%%%%%%%
% Seismica submission manuscript -- integrated analysis structure, version 60
% Supporting detail and diagnostic figures are in supplement60.tex.
%%%%%%%%%%%%%%%%%%%%%%%%%%%%%%%%%%%%%%%%%%%%%%%%%%%%%%%%%%%%%%%%%%%%%%%%%%%%%%%

\documentclass[]{seismica}
\usepackage{tabularx}
\usepackage{placeins}

% Color and TODO command for manuscript annotations
\usepackage{xcolor}

\newcommand{\todo}[1]{%
  \par\noindent
  \colorbox{yellow!40}{%
    \parbox{0.96\linewidth}{%
      \textbf{\textcolor{red}{TODO:}} #1
    }%
  }%
  \par
}

\newcommand{\checkme}[1]{\textcolor{red}{\textbf{[CHECK: #1]}}}

% Compile from the repository latex/ directory.
% Current external outputs are searched before bundled legacy assets.
\graphicspath{
  {/Users/thompsong/Library/CloudStorage/Box-Box/thompsong/3_Project_Documents/NASAprojects/201602_Rocket_Seismology/writing/falcon9_fireball_paper/outputs/000_correct_bchh_instrument_response/}
  {/Users/thompsong/Library/CloudStorage/Box-Box/thompsong/3_Project_Documents/NASAprojects/201602_Rocket_Seismology/writing/falcon9_fireball_paper/outputs/010_prepare_analysis_inputs/}
  {/Users/thompsong/Library/CloudStorage/Box-Box/thompsong/3_Project_Documents/NASAprojects/201602_Rocket_Seismology/writing/falcon9_fireball_paper/outputs/020_analyze_ksc_weather/}
  {/Users/thompsong/Library/CloudStorage/Box-Box/thompsong/3_Project_Documents/NASAprojects/201602_Rocket_Seismology/writing/falcon9_fireball_paper/outputs/030_reconcile_manual_event_catalogues/}
  {/Users/thompsong/Library/CloudStorage/Box-Box/thompsong/3_Project_Documents/NASAprojects/201602_Rocket_Seismology/writing/falcon9_fireball_paper/outputs/040_validate_event_catalogue_with_array_processing/}
  {/Users/thompsong/Library/CloudStorage/Box-Box/thompsong/3_Project_Documents/NASAprojects/201602_Rocket_Seismology/writing/falcon9_fireball_paper/outputs/040_validate_event_catalogue_with_array_processing/figures/}
  {/Users/thompsong/Library/CloudStorage/Box-Box/thompsong/3_Project_Documents/NASAprojects/201602_Rocket_Seismology/writing/falcon9_fireball_paper/outputs/050_prepare_event_waveform_products/}
  {/Users/thompsong/Library/CloudStorage/Box-Box/thompsong/3_Project_Documents/NASAprojects/201602_Rocket_Seismology/writing/falcon9_fireball_paper/outputs/050_prepare_event_waveform_products/figures/}
  {/Users/thompsong/Library/CloudStorage/Box-Box/thompsong/3_Project_Documents/NASAprojects/201602_Rocket_Seismology/writing/falcon9_fireball_paper/outputs/060_analyze_planar_array_propagation/}
  {/Users/thompsong/Library/CloudStorage/Box-Box/thompsong/3_Project_Documents/NASAprojects/201602_Rocket_Seismology/writing/falcon9_fireball_paper/outputs/070_measure_event_amplitudes_and_acoustic_seismic_coupling/}
  {/Users/thompsong/Library/CloudStorage/Box-Box/thompsong/3_Project_Documents/NASAprojects/201602_Rocket_Seismology/writing/falcon9_fireball_paper/outputs/070_measure_event_amplitudes_and_acoustic_seismic_coupling/figures/}
  {/Users/thompsong/Library/CloudStorage/Box-Box/thompsong/3_Project_Documents/NASAprojects/201602_Rocket_Seismology/writing/falcon9_fireball_paper/outputs/080_analyze_named_events_and_timing/}
  {/Users/thompsong/Library/CloudStorage/Box-Box/thompsong/3_Project_Documents/NASAprojects/201602_Rocket_Seismology/writing/falcon9_fireball_paper/outputs/090_estimate_acoustic_energetics/}
  {/Users/thompsong/Library/CloudStorage/Box-Box/thompsong/3_Project_Documents/NASAprojects/201602_Rocket_Seismology/writing/falcon9_fireball_paper/outputs/100_analyze_seismic_polarization/}
  {/Users/thompsong/Library/CloudStorage/Box-Box/thompsong/3_Project_Documents/NASAprojects/201602_Rocket_Seismology/writing/falcon9_fireball_paper/outputs/110_search_for_direct_seismic_waves/}
  {/Users/thompsong/Library/CloudStorage/Box-Box/thompsong/3_Project_Documents/NASAprojects/201602_Rocket_Seismology/writing/falcon9_fireball_paper/outputs/120_analyze_regional_detectability/}
  {/Users/thompsong/Library/CloudStorage/Box-Box/thompsong/3_Project_Documents/NASAprojects/201602_Rocket_Seismology/writing/falcon9_fireball_paper/outputs/130_generate_supplement_event_table/}
  {/Users/thompsong/Library/CloudStorage/Box-Box/thompsong/3_Project_Documents/NASAprojects/201602_Rocket_Seismology/writing/falcon9_fireball_paper/outputs/140_generate_publication_figures/}
  {/Users/thompsong/Library/CloudStorage/Box-Box/thompsong/3_Project_Documents/NASAprojects/201602_Rocket_Seismology/writing/falcon9_fireball_paper/outputs/150_assemble_zenodo_release/}
  {figures/}
}

% These boxes keep the manuscript compilable while the two remaining main
% figures are regenerated. They should not appear in a submitted PDF.
\newcommand{\pendingfigure}[1]{%
  \fbox{\parbox[c][55mm][c]{0.92\textwidth}{\centering\textbf{Pending figure}\par#1}}}

\title{Close-Range Seismo-Acoustic Observations of the 2016 Falcon 9 AMOS-6 Failure at Cape Canaveral}
\shorttitle{Close-range observations of the 2016 Falcon 9 explosion}

% Confirm department-level present affiliations for Farrell and Mehta before
% submission.
\author[1]{Glenn Thompson\thanks{Corresponding author: \texttt{thompsong@usf.edu}; ORCID: 0000-0002-9173-0097}}
\author[2,3]{Robert G. Brown}
\author[1,4]{Alexandra Farrell}
\author[2,3]{John Kiriazes}
\author[1]{Stephen R. McNutt}
\author[1,5]{Christopher Mehta}
\affil[1]{School of Geosciences, University of South Florida, Tampa, Florida, USA}
\affil[2]{NASA Kennedy Space Center, Florida, USA}
\affil[3]{Present address: Florida Space Institute, University of Central Florida, Orlando, Florida, USA}
\affil[4]{Present address: University of Alaska Fairbanks, Fairbanks, Alaska, USA}
\affil[5]{Present address: Howard University, Washington, DC, USA}

\credit{Conceptualization}{Glenn Thompson}
\credit{Methodology}{Glenn Thompson}
\credit{Software}{Glenn Thompson}
\credit{Validation}{Glenn Thompson}
\credit{Formal Analysis}{Glenn Thompson}
\credit{Investigation}{Glenn Thompson, Robert G. Brown, Alexandra Farrell,
John Kiriazes, Stephen R. McNutt, Christopher Mehta}
\credit{Resources}{Robert G. Brown, John Kiriazes}
\credit{Data Curation}{Glenn Thompson}
\credit{Visualization}{Glenn Thompson}
\credit{Project Administration}{Glenn Thompson, Robert G. Brown}
\credit{Writing - Original draft}{Glenn Thompson}
\credit{Writing - Review \& Editing}{Glenn Thompson, Robert G. Brown,
John Kiriazes, Stephen R. McNutt}

\begin{document}

\makeseistitle{
  \begin{summary}{Abstract}
  On 1 September 2016, a Falcon~9 and its AMOS-6 payload were destroyed
  during propellant loading before a planned static-fire test at Space Launch
  Complex~40, Cape Canaveral. A broadband seismometer and three-element
  infrasound array 1.4~km away recorded the complete sequence. A gap-free
  48-h review found no seismo-acoustic precursor during the 37~h before the
  first upper-stage event. We then compared relative source-event timing from
  the public USLaunchReport video with independently picked BCHH pressure
  arrivals. The principal explosion occurred about 3.57--3.64~s after the
  first event in that recording, whereas its DD2--DD3 pressure arrivals
  followed by only 3.501~s. The resulting 0.07--0.14-s differential travel-time
  reduction is independent of absolute camera synchronization. Two later
  payload-associated events give smaller reductions of about 0.034 and
  0.054~s. Higher-frame-rate proprietary SpaceX recordings, including the
  UCS-3 camera near BCHH and cameras at SLC-40, independently support the event
  identifications and frame-scale timing. A simple weak-shock sensitivity
  model predicts the same ordering and comparable differential advances. At BCHH, however, the
  preferred local array-crossing speed of the principal blast was only
  352.5~m~s$^{-1}$ and the measured 1.39-kPa positive overpressure corresponds
  to $M\approx1.006$, consistent with a blast that was more nonlinear closer
  to the vehicle and had decayed to a very weak shock by 1.4~km. The full
  catalogue contains 153 accepted pressure transients over 26~min. The
  principal explosion produced a peak level of about 157~dB re 20~$\mu$Pa and
  a hemispherical-equivalent acoustic-energy estimate of $2.2\times10^9$~J.
  Pressure and three-component peak ground velocity were nearly proportional,
  with a median seismic/acoustic coupling coefficient of
  $2.3\times10^{-6}$~(m~s$^{-1}$)~Pa$^{-1}$.
  \end{summary}
  %
  \begin{summary}{Non-technical summary}
  A seismic and low-frequency sound station approximately 1.4~km from the
  launch pad recorded the 2016 Falcon~9 accident at Cape Canaveral. We found
  no comparable signal in the 37~h before the first event. The accident was
  not a single explosion: 153 accepted multichannel pressure transients
  occurred over approximately 26~min. For the first few seconds, the public
  USLaunchReport video provides an independent source-event clock. We use only
  time differences measured within that recording, so the camera does not need
  to be synchronized to the geophysical station. The strongest explosion
  occurred about 3.57--3.64~s after the first visible event but reached the
  BCHH microphones only 3.501~s after the first pressure arrival. Proprietary
  SpaceX videos provide higher-frame-rate checks on these event identifications.
  The difference indicates that the stronger blast propagated faster over part
  of the path.
  Its speed across the small BCHH array was nevertheless close to ordinary
  sound speed, showing that the disturbance had weakened substantially by
  1.4~km range. The observations provide an independently timed forensic record
  of the failure when direct site access and engineering telemetry are not
  available to outside investigators.
  \end{summary}
}

\section{Introduction}
\label{sec:introduction}

Rocket launches generate strong atmospheric and ground-motion signals, but
catastrophic launch failures are rarely recorded at close range by calibrated
seismic and infrasound instruments. Most seismo-acoustic studies of major
technological explosions have instead used regional observations, for which
atmospheric ducting and attenuation strongly influence arrival times and
amplitudes \citep{ottemoller_2008,green_2008,ceranna_2009,smink_2019,
pilger_beirut_2021,kim_pasyanos_2022}. Rocket-focused studies demonstrate
complementary capabilities: routine launches are detectable infrasonically at
large distances \citep{pilger_rocket_2021}, dense seismic deployments can
resolve launch-generated body and surface waves
\citep{schuster_rocketquake_2025}, and regional infrasound can locate and
estimate the effective yield of a large static-fire explosion
\citep{silber_scamfer_newglenn_2026}.

On 1 September 2016, a Falcon~9 was destroyed during propellant loading in the
countdown for a planned static-fire test at Space Launch Complex~40 (SLC-40),
Cape Canaveral Air Force Station (now Cape Canaveral Space Force Station),
Florida. The AMOS-6 payload was lost and the launch complex was extensively
damaged. The only publicly available original recording is the 5-min-40-s raw
USLaunchReport video and camera audio \citep{uslaunchreport_amos6_video_2016};
an edited analysis by \citet{manley_spacex_explosion_2016} was used only for
context. 

The University of South Florida (USF) BCHH station, approximately 1.4~km from
SLC-40 on adjacent Kennedy Space Center property, fortuitously recorded the
complete sequence. The initial forensic analysis was documented in an
unpublished USF technical report prepared for SpaceX in November~2016
\citep{thompson_mcnutt_fireball_2016} and subsequently presented publicly at
the 2017 AGU Fall Meeting \citep{thompson_agu_falcon9_2017}, before the
present reanalysis. The geophysical chronology was developed independently of
the later engineering root-cause determination.

The central advantage of this dataset is the combination of continuous
recording, empirical calibration, array geometry, and proximity. It provides
subsecond timing and local pressure--ground comparisons that regional records
rarely preserve, while public imagery supplies an independently inspectable
record of selected stages. Its limits are equally clear: one compact array can
constrain direction and apparent speed, but not source range or the origin of
every later transient.

Here we first test for a detectable energetic precursor and establish the
initial event that anchors the chronology. We then compare event-to-event
intervals measured independently in a single SpaceX video with the corresponding
pressure-arrival intervals at BCHH. This differential comparison avoids the
circularity of using source times obtained from the same pressure arrivals to
test propagation speed. Array processing then constrains the wavefront locally
at BCHH, while pressure, ground motion, coupling, and acoustic-energy estimates
characterize the broader sequence. A separate regional search tests whether the
principal explosion can be associated with publicly archived records at much
greater distances. The broader objective is to show what a modest open
geophysical deployment can establish independently, and what still requires
engineering telemetry, denser spatial sampling, or detailed atmospheric
modelling. We use \emph{close range}, rather than the strict acoustic term
\emph{near field}: a 1.4-km wavelength at 350~m~s$^{-1}$ corresponds to
approximately 0.25~Hz, below the principal analysis band. Supporting methods,
diagnostics, frame picks, and sensitivity results are provided in the
Supplementary Material.


\section{Data and general processing}
\label{sec:data}

\subsection{BCHH array and geometry}

BCHH comprised a three-component Nanometrics Trillium Compact Posthole
broadband (>0.083 Hz) seismometer near the centre of a three-element infraBSU microphone (>0.048 Hz) array with an
aperture of approximately 30~m (Figure~\ref{fig:site_array}). All six channels
were recorded continuously at 250~Hz by a GPS-synchronized Nanometrics
Centaur digitizer. The seismometer was buried approximately 0.5~m; the infrasound sensors
were mounted ~30~cm above ground on garden stakes in shrubs to suppress wind noise.

On 6 October 2016, as Hurricane Matthew approached KSC, co-author RGB 
removed the three microphones and stakes, but the cables remained on 
the ground. The cable ends were subsequently surveyed with precise GPS on
2 November 2016. Position errors of approximately 1--2~m are therefore plausible.
Source distances ranged from 1408.3 to 1440.0~m
(Supplementary Table~S1). DD2 was used as the reference
pressure channel, as it had the best signal-to-noise ratio.

The USLaunchReport video camera was located on the roof of a building in the Fire Training Area, 
providing a clear view of the launch complex. Verification of the camera's line of sight is shown in Supplementary Figure~S1.

\begin{figure*}[!tbp]
  \centering
  \includegraphics[width=0.96\textwidth]{fig01_slc40_bchh_location.pdf}
  \caption{Locations used in the study (A) Position of the BCHH seismo-acoustic array
   and the USLaunchReport camera. The launch complex was SLC-40 (yellow triangle).
  (B) Three-component seismometer and three-element infrasound
  array with an aperture of approximately 30~m.}
  \label{fig:site_array}
\end{figure*}
\FloatBarrier

\subsection{Calibration and waveform preparation}

The seismic response was removed using the nominal calibration for the
Trillium Compact Posthole--Centaur combination. 1-m lift tests (Supplementary Section~S1) were subsequently performed and 
determined sensitivities of 8.10, 0.69, and
10.0~counts~Pa$^{-1}$, respectively for infraBSU sensor channels DD1, DD2, and DD3 (Supplementary Table~S1).
These empirical factors \citep{thompson_mcnutt_fireball_2016} are retained here (Supplementary Table~S1).
Results from other launches reported here are consistent with these empirical factors.
The $\pm20\%$ pressure uncertainty reflects the imprecision and repeatability of the empirical lift test procedure. 
Manual arrival uncertainty was approximately 0.01~s for the clearest impulses and
greater for weak, emergent, or overlapping signals.

Response removal used a 0.02, 0.05, 80, and 100~Hz four-corner prefilter and a
60-dB water level. A centred 1.0-s moving median was then removed independently
from each pressure channel to suppress slow baseline variation without an
additional physical-amplitude high-pass filter. Seismic traces remained in
ground velocity and pressure traces in Pa.

\subsection{Video, meteorological, and regional data}

The USLaunchReport footage \citep{uslaunchreport_amos6_video_2016} provided the public visual and camera-audio record and is the primary video source for relative timing of the four opening events. Event intervals were measured from video presentation timestamps (PTS) within this single recording; its absolute clock need not be synchronized to BCHH. Four proprietary videos obtained from SpaceX were used as independent higher-quality checks on event identification and timing. These include the 60-fps UCS-3 remote camera at 28.575108$^\circ$N, 80.573638$^\circ$W, about 155~m from BCHH and between SLC-40 and SLC-41, and cameras at SLC-40. The UCS-3 burned-in timecode advances at 30 frames~s$^{-1}$, so each displayed timecode value occurs on two successive video images. No event interval is formed by mixing timestamps from different cameras. Data were sought from an infrasound array at CCAFS operated by the Air Force Technical Applications Center (AFTAC) but were ultimately inaccessible. Regional pressure records to 1800~km
and vertical seismic records to 600~km were obtained through EarthScope FDSN
services; IU.DWPF, 97.7~km from SLC-40, was the nearest permanent station
examined \citep{asl_gsn_iu_1988}.

Kennedy Space Center meteorological observations were examined at two spatial scales. Weather-tower measurements supplied temperature, relative humidity, and near-surface winds for the close-range SLC-40--BCHH propagation correction. To characterize the vertically extended atmosphere relevant to regional infrasound propagation, we additionally examined KSC 915-MHz and 48-MHz radar wind-profiler observations and a 12:00~UTC low-resolution AMPS radiosonde sounding. The profilers constrained winds through the lower and middle atmosphere, including 48-MHz profiles at 13:04 and 13:09~UTC bracketing the principal explosion, while the AMPS sounding supplied temperature and winds into the lower stratosphere. No atmospheric ray tracing was performed. Additional details are given in Supplementary Section~S2.

The tower observations were used only for this close-range timing reference.
Vertically extended KSC 915-MHz and 48-MHz wind-profiler observations and the
AMPS radiosonde sounding were treated separately as regional-propagation
diagnostics and were not used to alter the close-range reduction.

\subsection{Analysis workflow and timing coordinates}

The analysis was implemented in Python as a reproducible sequence of Jupyter
notebooks \citep{kluyver2016jupyter}. Shared configuration and intermediate
products supplied common geometry, receiver arrival picks, event metadata,
propagation corrections, and analysis parameters. The reviewed BCHH arrival
picks, public-video PTS picks, and proprietary SpaceX frame observations are
stored as separate metadata products so that measured receiver timing, independent video timing, and
model-derived source-time coordinates are not conflated. Additional
implementation and sensitivity details are in Supplementary Sections~S1--S10.

Three time coordinates were kept distinct: observed receiver UTC,
DD2-referenced reduced arrival time, and model-derived source time. For sensor
$i$, the DD2-referenced reduced arrival time was
\begin{equation}
 t_{i,\mathrm{red}}=t_i-\frac{r_i-r_{\mathrm{DD2}}}{c_{\mathrm{eff}}},
 \label{eq:reduced_time}
\end{equation}
where $t_i$ is observed receiver time, $r_i$ the source-to-sensor distance,
$r_{\mathrm{DD2}}$ the nominal source-to-DD2 distance, and
$c_{\mathrm{eff}}$ the meteorologically derived effective speed. Catalogue
times therefore represent arrivals reduced to DD2. Legacy source-time
coordinates subtract the complete source-to-DD2 travel time as well; because
those source times are derived from the BCHH arrivals, they are not used below
as independent evidence for propagation speed.

\FloatBarrier


\section{Forensic chronology and precursor search}
\label{sec:precursor}

A gap-free 48-h interval extending from 37~h before to 11~h after the first
upper-stage signal was manually scanned using Antelope's \texttt{dbpick}
program (version~5.6) \citep{brtt_antelope}. All three pressure and all three
seismic channels were viewed together. This review served two purposes: to
look explicitly for a conspicuous precursor before the first accepted event,
and to pick impulsive signals during the accident sequence.

SpaceX were particularly attentive to any precursor signals. None was
identified during the 37~h before the initial upper-stage explosion, and no
additional explosion-related impulse was found in the remainder of the 48-h
review window, which extends about 10.5~h beyond the last prominent activity
phase. The six-channel screen constrains signals comparable to the accepted
array-wide impulses under the observed background, but cannot exclude weak,
spatially local, non-impulsive, internal, or poorly coupled processes. The
result is therefore a negative geophysical observation, not a determination
of the internal failure mechanism.


\section{Opening sequence and propagation timing}
\label{sec:opening_timing}

\subsection{Close-range acoustic reference}

For the close-range acoustic travel-time reference, still-air sound speed was
calculated from the contemporaneous KSC weather-tower temperature and
relative-humidity measurements described above. Wind vectors from the KSC weather-tower observations were
projected onto the SLC-40--BCHH propagation direction. To characterize the wind
over the nominal 0--65-m source-height interval, the available tower wind
observations within that height range were combined as a vector mean. Because
the lowest valid wind observations were near 40~m, the lowest measured wind
vector was extended downward at a constant value to estimate the nominal
0--65-m path average; the surface layer was therefore not directly sampled.
The path-parallel wind component was added to the still-air sound speed to
obtain the nominal effective acoustic speed, 351~m~s$^{-1}$, used for
reference travel times and catalogue reduction.

Seven 5-min weather-tower epochs from 13:05 to 13:35~UTC gave a mean
temperature of 26.4$\pm0.3\,^{\circ}$C and mean relative humidity of
83.9$\pm3.7\%$. The tower winds varied little, with an overall vector mean of
5.4~m~s$^{-1}$ from 144$^\circ$. Averaging the available winds over the
nominal 0--65-m source-height interval gave a path-representative vector mean
of 4.6~m~s$^{-1}$ from 145$^\circ$, corresponding to a 2.6~m~s$^{-1}$
tailwind component toward BCHH. Together with a still-air sound speed of
348.3~m~s$^{-1}$, this gives the nominal effective acoustic speed of
351.0~m~s$^{-1}$. The corresponding travel times are about 4.01~s to DD2
(1408.3~m) and 4.05~s to the central seismometer (1419.9~m).

\subsection{Independent video and BCHH event timing}

The public USLaunchReport recording \citep{uslaunchreport_amos6_video_2016}
provides the primary source-event timing because it is publicly accessible and
therefore independently reproducible. Its absolute clock is not synchronized to
BCHH, so only event-to-event intervals within that same recording are used. The
first visible upper-stage flash, at PTS 71.7675~s, is the reference event. The
first unmistakable frame of the principal explosion occurs at PTS 75.4048~s.
Inspection of the immediately preceding frames indicates that the visual onset
probably began no more than one or two frames earlier. At 29.97~frames~s$^{-1}$
this gives a preferred principal source interval of 3.5706--3.6373~s after the
initial event. The preferred payload-event picks are PTS 84.2813 and
84.8819~s, giving intervals of 12.5138 and 13.1144~s, respectively.

Four proprietary SpaceX recordings provide independent higher-quality checks
on event identification and timing, including the 60-fps UCS-3 remote camera
near BCHH/Astronaut Beach House and cameras at SLC-40. Their clocks are not
assumed to be synchronized with one another or with BCHH, and no timestamps
are mixed between cameras. Camera geometry, timecode handling, and the
separate SpaceX timing measurements are documented in Supplementary
Section~S4.

DD2 and DD3 arrivals for the upper-stage event, principal explosion, payload
impact, and payload explosion were manually reviewed and stored as direct
receiver picks. For a later event $j$, the observed same-channel arrival
interval is
\begin{equation}
 \Delta t_{\mathrm{BCHH},j}=t_{j,\mathrm{arr}}-t_{1,\mathrm{arr}},
\end{equation}
and the independent video source interval is
$\Delta t_{\mathrm{video},j}$. For sources at the same nominal location, a
constant 351~m~s$^{-1}$ travel time cancels from the interval comparison. The
nominal arrival predicted from the first observed arrival is simply the first
arrival plus $\Delta t_{\mathrm{video},j}$; a source-location correction
$(r_j-r_1)/351$ can be added when justified. We define
\begin{equation}
 \delta t_j=\Delta t_{\mathrm{video},j}-\Delta t_{\mathrm{BCHH},j},
 \label{eq:differential_advance}
\end{equation}
so positive $\delta t_j$ means that event $j$ required less propagation time
to reach BCHH than the initial event. This differential construction is the
main timing test; the 351~m~s$^{-1}$-reduced source times are retained only as
a workflow coordinate.

The median DD2--DD3 arrival interval from the first event to the principal
explosion is 3.501014~s. The preferred USLaunchReport source interval of
3.5706--3.6373~s therefore gives a differential propagation-time reduction of
0.0696--0.1363~s. This comparison does not use the
351~m~s$^{-1}$-reduced principal source time and is therefore independent of
that reduction.

For the payload impact, the public video gives a source interval of 12.5138~s,
compared with a median BCHH interval of 12.480232~s, giving
$\delta t=0.0336$~s. The payload-associated explosion occurs 13.1144~s after
the initial event in the public video and 13.060170~s after the initial BCHH
arrival, giving $\delta t=0.0542$~s. These values are based on a 29.97-fps
recording and should be treated as frame-scale estimates rather than
millisecond-accurate source times. The payload source positions are also less
certain than those of the stage explosions, so source-range differences matter
for residuals of a few hundredths of a second.

The proprietary SpaceX videos independently reproduce the same qualitative
timing pattern. UCS-3 gives a principal-event differential reduction of about
0.10--0.17~s and smaller reductions of about 0.020 and 0.023~s for the two
payload events; the SLC-40 views further support the event identifications and
onset bounds. The underlying frame observations and within-camera intervals
are given in Supplementary Section~S4.

\begin{figure*}[!tbp]
  \centering
  \includegraphics[width=0.98\textwidth]{fig04_multimodal_alignment_upper_stage.pdf}
  \caption{Video, camera audio, pressure, and ground-velocity alignment for
  the first upper-stage event. Camera audio is corrected for its estimated
  source-to-camera travel time and BCHH traces for sensor-specific reference
  travel times. This figure establishes the visual identity of the first event;
  it is not used to assign absolute UTC to subsequent video events. Video
  frames copyright 2016 USLaunchReport.com/Michael Wagner; reproduced with
  permission \citep{uslaunchreport_amos6_video_2016}.}
  \label{fig:multimodal}
\end{figure*}

\begin{table*}[!tbp]
\caption{Opening-sequence timing and pressure. Public-video intervals are
measured within the USLaunchReport recording relative to PTS 71.7675~s; no
absolute camera--BCHH synchronization is assumed. BCHH intervals are the
median of same-channel DD2 and DD3 arrival differences. The differential
advance is $\Delta t_{\rm video}-\Delta t_{\rm BCHH}$, so positive values mean
the later signal required less propagation time than the initial event. The
principal interval spans the preferred 0--2-frame onset range. Proprietary
SpaceX timing checks are given in Supplementary Section~S4.}
\label{tab:principal_events}
\centering
\small
\begin{tabularx}{\textwidth}{Xccccc}
\hline
Signal &
\shortstack{USLaunchReport interval\\from initial (s)} &
\shortstack{Median BCHH interval\\from initial (s)} &
\shortstack{Differential\\advance (s)} &
$p_+$ (Pa) &
\shortstack{$L_{\mathrm{peak}}$\\(dB)} \\
\hline
Upper-stage event & 0 & 0 & reference & 38.8 & 125.8 \\
Principal explosion & 3.571--3.637 & 3.501014 & 0.070--0.136 & 1390 & 156.8 \\
Payload impact & 12.5138 & 12.480232 & 0.0336 & 279.1 & 142.9 \\
Payload-associated explosion & 13.1144 & 13.060170 & 0.0542 & 297.0 & 143.4 \\
\hline
\end{tabularx}
\end{table*}

\subsection{Weak-shock timing sensitivity}

The positive overpressure of each named event was also used in a deliberately
simple weak-shock sensitivity calculation. Local Mach number followed the
normal-shock pressure-jump relation. To test whether faster propagation close
to the source could be compatible with a nearly acoustic wave at BCHH, the
measured pressure was extrapolated inward with $\Delta p\propto r^{-1}$ to a
10-m cutoff and travel time integrated along the path. This is a sensitivity
model, not a source-pressure inversion; the $1/r$ law is not validated in the
near-source region. The quantity compared with Equation~\ref{eq:differential_advance}
was the modeled propagation advance relative to the initial event, rather than
an absolute source time.

The pressure dependence of these timing differences is consistent with the
weak-shock sensitivity calculation. Using the measured positive pressures and
the simple inward $1/r$ model, the modeled propagation advance relative to the
initial event is approximately 0.095~s for the principal explosion,
0.019~s for the payload impact, and 0.021~s for the payload explosion. The
principal prediction lies within the preferred public-video range. The two
public-video payload residuals are somewhat larger than the simple model, while
the proprietary UCS-3 estimates (about 0.020 and 0.023~s) are closer. Given the
frame scale and uncertain payload source locations, we regard the common
ordering and order of magnitude as more significant than the event-by-event
numerical agreement. Figure~\ref{fig:differential_timing} shows the underlying
camera-predicted and reviewed pressure arrivals; the weak-shock values above
remain a sensitivity calculation rather than a fit to those traces.

\begin{figure*}[!tbp]
  \centering
  \IfFileExists{/Users/thompsong/Library/CloudStorage/Box-Box/thompsong/3_Project_Documents/NASAprojects/201602_Rocket_Seismology/writing/falcon9_fireball_paper/outputs/140_generate_publication_figures/fig_named_event_video_predicted_vs_observed_arrivals.pdf}{%
    \includegraphics[width=0.98\textwidth]{fig_named_event_video_predicted_vs_observed_arrivals.pdf}%
  }{%
    \pendingfigure{Camera-predicted and reviewed BCHH pressure arrivals for the
    principal explosion, payload impact, and payload explosion.}%
  }
  \caption{Camera-predicted and reviewed BCHH pressure-arrival times for the
  three later named events. Each row shows a reviewed USLaunchReport frame
  (left) and baseline-corrected DD2 and DD3 pressure traces (middle and right),
  with time measured from the reviewed upper-stage arrival on that channel.
  Solid blue and green lines mark the midpoints of the USLaunchReport and UCS-3
  arrival predictions; horizontal double-headed arrows show their timing
  brackets where an onset range is available. Red dashed lines mark the
  reviewed pressure-arrival picks. The principal-event camera brackets
  overlap; payload timings are preferred point picks. All predictions use
  within-camera event intervals and do not assume absolute camera--BCHH clock
  synchronization. The payload predictions assume an unchanged acoustic path,
  so source motion or position changes can also contribute to residuals.
  Video frames copyright 2016 USLaunchReport.com/Michael Wagner; reproduced
  with permission \citep{uslaunchreport_amos6_video_2016}.}
  \label{fig:differential_timing}
\end{figure*}
\FloatBarrier


\section{Event catalogue and temporal evolution}
\label{sec:temporal_evolution}

Conspicuous impulsive pressure signals readily visible across the microphone
array were manually picked, usually giving one candidate arrival on each
microphone. The rest of the catalogue pipeline was automated. For association,
each candidate arrival was reduced to the DD2 reference by subtracting the
predicted differential travel time from SLC-40 to that sensor using
351~m~s$^{-1}$. A moving 0.04-s window, the smallest coincidence window on a
stable association-count plateau, was then applied to identify coincident
arrivals across the array and group them into events. This simple time-based
event association was performed using GISMO \citep{thompson_gismo_2017}.

Catalogue membership was separated from subsequent waveform-quality
classification. Associated multichannel pressure transients constituted the
event catalogue; the catalogue therefore records accepted multichannel
pressure transients rather than independently verified physical explosions.
DD1 was sporadically noisier and, overall, correlated less consistently with
the other two pressure channels, so DD2 and DD3 provided the most stable
waveform-quality pair. The usable-waveform subset required robust peak
SNR $\geq3$ and DD2--DD3 correlation $\geq0.5$, whereas the core subset
required SNR $\geq5$ and correlation $\geq0.7$. These waveform-quality labels
did not determine catalogue membership.

The final catalogue contains 153 accepted multichannel pressure transients.
These occur in four descriptive phases containing 39, 4, 90, and 6 entries,
respectively; 14 entries fall outside those intervals
(Figure~\ref{fig:overview}). Phase~I spans the first 70~s of activity;
Phases~II--IV mark later visually separated clusters. They are not inferred
source regimes. Of the 153 entries, 138 met usable-waveform and 107 met core
criteria. The independently defined planar classes contained 51 high-quality,
45 intermediate, and 57 weak or broad solutions. Median robust peak SNR was
7.86, median DD2--DD3 correlation 0.876, and median corrected span among
associated picks 4.87~ms. These nested counts serve different analytical
purposes and should not be interpreted as counts of independent explosions.

A prolonged signal lasting approximately 3.2~min followed Phase~I; it is
visible on DD2 pressure as well as the seismic channels, and the public footage
shows vigorous fire across SLC-40 during this interval.

\begin{figure*}[!tbp]
  \centering
  \includegraphics[width=0.98\textwidth]{fig03_bchh_30min_overview_zp_true.pdf}
  \caption{Thirty-minute six-channel BCHH raw-data overview (no time reduction).
  The upper traces are pressure and the lower traces are three-component ground
  velocity. All traces use a zero-phase four-pole 0.5--25-Hz display filter and
  are independently scaled by their robust 99.3rd-percentile amplitude;
  amplitudes therefore cannot be compared between rows. The shaded Phase-I--IV
  intervals contain 39, 4, 90, and 6 catalogue transients, respectively; 14
  transients fall outside those intervals. The prolonged post-Phase-I signal is
  marked descriptively rather than as a source classification.}
  \label{fig:overview}
\end{figure*}

The phase structure is also evident in continuous energetic diagnostics from
the same records (Figure~\ref{fig:sequence_energy_rates}). Five-second
increments of high-pass-filtered DD2 acoustic energy emphasize the principal
explosion and later bursts, whereas the three-component ground-motion energy
proxy separates Phases~II--IV particularly clearly. The proxy is the
band-limited increment of $\int(v_Z^2+v_R^2+v_T^2)\,dt$ and is not interpreted
as radiated seismic energy.

\begin{figure*}[!tbp]
  \centering
  \IfFileExists{/Users/thompsong/Library/CloudStorage/Box-Box/thompsong/3_Project_Documents/NASAprojects/201602_Rocket_Seismology/writing/falcon9_fireball_paper/outputs/090_estimate_acoustic_energetics/full_sequence_cumulative_acoustic_and_ground_motion.pdf}{%
    \includegraphics[width=0.98\textwidth]{full_sequence_cumulative_acoustic_and_ground_motion.pdf}%
  }{%
    \pendingfigure{Continuous sequence acoustic-energy and ground-motion-proxy
    rates from Notebook 090.}%
  }
  \caption{Continuous energetic evolution of the 26-min accident sequence in
  5-s bins. (a) Hemispherical acoustic-energy increments from DD2 after a
  zero-phase 0.5-Hz high-pass filter. (b) Three-component ground-motion energy
  proxy after 0.5--20-Hz filtering. The latter is a local measure of recorded
  ground motion, not radiated seismic energy. Shading marks the four descriptive
  phase windows and hatching the prolonged post-Phase-I signal.}
  \label{fig:sequence_energy_rates}
\end{figure*}
\FloatBarrier


\section{Waveform morphology and array propagation}
\label{sec:array_propagation}

For the high-quality planar subset, pressure and vertical-ground-velocity
waveforms were aligned to the event reference and normalized independently
before calculating pointwise medians and percentile envelopes. The waveform
ensemble is used to describe repeatable morphology, not to infer source timing
or Mach number.

The 51 high-quality planar events show broadly repeatable impulsive compression
followed by a smaller and longer rarefaction, whereas vertical ground velocity
has a substantially longer oscillatory coda
(Figure~\ref{fig:waveform_ensemble}). The strongest pressure pulses are
Friedlander-like in form; classical symmetric N-wave morphology was not
observed. Many weaker pressure impulses approach the seismic and acoustic noise
floors. Waveform shape is therefore consistent with blast-wave propagation but
is not by itself a Mach-number measurement.

\begin{figure*}[!tbp]
  \centering
  \includegraphics[width=0.88\textwidth]{high_quality_normalized_pressure_and_dhz_waveforms.pdf}
  \caption{Independently normalized pressure and vertical-ground-velocity
  waveforms for the 51 events in the high-quality planar class. Pale curves
  are individual events, the dark curve is the pointwise median, and shading
  spans the 16th--84th percentiles.}
  \label{fig:waveform_ensemble}
\end{figure*}

A time-domain planar-wave grid search estimated back azimuth and apparent
velocity by maximizing multichannel waveform agreement. All catalogue entries
were processed, while physical interpretation emphasized quality-controlled
solutions. A fixed-source circular-wavefront calculation tested sensitivity to
the 1.4-km source range; because the three-element array provides only two
independent relative delays, it was not expected to resolve both source range
and propagation speed.

Time-domain propagation quality was classified independently. The high-quality
planar class required maximum coherence $\geq0.90$, back-azimuth standard
deviation $\leq0.5^\circ$, and apparent-speed standard deviation
$\leq3.0$~m~s$^{-1}$ across nine perturbations of the event scoring window.
Intermediate solutions used corresponding thresholds of 0.75, 1.0$^\circ$,
and 8.0~m~s$^{-1}$; remaining solutions were classified as weak or broad.

Independent frequency-domain Bartlett beamforming was performed using InfraPy
version~0.4.1.0, developed at Los Alamos National Laboratory (LANL)
\citep{webster_infrapy_2025}, with all three pressure channels. The primary
validation used 0.5-s windows centred on the catalogue times and a 2--20-Hz
band. If $C$ denotes multichannel coherence for $N=3$ sensors, the Bartlett
statistic was
\begin{equation}
 F=\frac{(N-1)C}{1-C}.
\end{equation}
It is a monotonic transformation of coherence, not a measure of pressure,
energy, or yield. An empirical screening threshold was the 99.5th percentile
of one nonoverlapping 0.5-s pre-event background window for each catalogue
entry. A separate scan of 17,839 overlapping windows, positioned without
reference to catalogue times, tested catalogue recovery and searched for
additional coherent activity.

Source-to-array integrated travel time and array-crossing apparent velocity are
different measurements. The timing comparison above tests integrated
propagation over the approximately 1.4-km path, whereas array processing
measures the wavefront locally across the approximately 30-m BCHH aperture and
is sensitive to metre-scale reconstructed-position errors.

In the event-centred InfraPy Bartlett analysis, median back azimuth was
200.7$^\circ$ compared with the geometric 199.4$^\circ$; 145 of 153 events
were within 5$^\circ$ of SLC-40. Median apparent velocity was
354.6~m~s$^{-1}$, and 113 events exceeded the empirical background threshold
$F=3.02$. The blind scan matched 129 catalogue entries (84.3\%) one-to-one;
recovery increased to 75 of 77 entries with event-centred $F\geq5$ and 41 of
42 with $F\geq10$.

The independent time-domain analysis gave a high-quality coherence-weighted
direction of 200.7$^\circ$ from projected grid north and an apparent velocity
of 352.0~m~s$^{-1}$. Conditional sensor-position tests show that 1--2-m
reconstruction errors across the aperture can exceed the observed direction
and velocity dispersions (Supplementary Section~S5). The finite-distance
calculation did not resolve range and slightly reduced coherence; the planar
result is therefore preferred.

For the principal explosion, the preferred time-domain array-crossing solution
was 352.5~m~s$^{-1}$ with coherence 0.9984. A separate short-window Bartlett
solution gave approximately 369.4~m~s$^{-1}$ with $F\simeq4.49$. The measured
1.39-kPa positive overpressure provides an independent constraint: interpreted
as a normal-shock pressure jump, it gives $M=1.0059$ at BCHH and a local speed
of about 353.0~m~s$^{-1}$ after adding the path-parallel tailwind. A local
shock speed of 369.4~m~s$^{-1}$ would require overpressure of order 13~kPa,
about an order of magnitude larger than measured. The higher Bartlett value is
therefore unlikely to be the true local shock speed.

\begin{figure*}[!tbp]
  \centering
  \IfFileExists{/Users/thompsong/Library/CloudStorage/Box-Box/thompsong/3_Project_Documents/NASAprojects/201602_Rocket_Seismology/writing/falcon9_fireball_paper/outputs/060_analyze_planar_array_propagation/combined_04b_08_array_summary.pdf}{%
    \includegraphics[width=0.98\textwidth]{combined_04b_08_array_summary.pdf}%
  }{%
    \pendingfigure{Combined time-domain planar and event-centred Bartlett
    array summary from the streamlined propagation notebooks.}%
  }
  \caption{Array-processing results through the 153-transient sequence.
  Time-domain planar back azimuth and apparent velocity are shown with the
  independent event-centred Bartlett statistic. Dashed lines mark the geometric
  SLC-40 direction, meteorologically predicted effective acoustic speed, and
  empirical Bartlett background threshold. Bartlett $F$ measures coherent
  plane-wave fit and is not a pressure or energy measure.}
  \label{fig:array_summary}
\end{figure*}
\FloatBarrier


\section{Seismic/acoustic coupling}
\label{sec:coupling}

Catalogue and regression pressures were peak-to-peak amplitudes of the
differentially aligned median waveform; named-event pressures were medians of
independently measured channel extrema. Three-component vector peak ground
velocity (PGV) was measured over a pressure-defined seismic window. A log--log
model
\begin{equation}
 \log_{10}p_{\mathrm{p2p}}=a+b\log_{10}(\mathrm{PGV})
 \label{eq:power_law}
\end{equation}
was fitted to entries with pressure-stack and vector-PGV SNR both $\geq5$ and
event-window PGV at least 70\% of the maximum in the saved approximately
1.1-s segment. This criterion does not establish full event/coda capture.
Conditional percentile ranges were obtained by resampling overlapping
measurement groups together, with catalogue-entry resampling retained as a
comparison.

Forty-six catalogue entries occupying 45 overlapping-measurement groups met
the joint criteria for pressure--PGV regression. The fitted slope was 0.962
(overlap-group bootstrap 2.5--97.5 percentiles 0.871--1.011), $R^2=0.953$,
and Spearman correlation 0.931 (Figure~\ref{fig:pressure_pgv}). Excluding all
selected rows in the principal measurement group gave $b=0.943$ and
$R^2=0.905$ for 44 entries. The median selected-entry pressure-to-vector-PGV
ratio was $4.34\times10^5$~Pa per (m~s$^{-1}$), equivalently
$2.31\times10^{-6}$~(m~s$^{-1}$)~Pa$^{-1}$.

\begin{figure*}[!tbp]
  \centering
  \includegraphics[width=0.82\textwidth]{publication_pressure_pgv_regression.pdf}
  \caption{Peak-to-peak pressure versus three-component vector PGV for 46
  measurement-quality catalogue entries. The solid curve is the fitted power
  law. The dashed line shows the median proportional relation,
  $p_{\mathrm{p2p}}=4.34\times10^5\,\mathrm{PGV}$, equivalently
  $\mathrm{PGV}=2.31\times10^{-6}p_{\mathrm{p2p}}$, in SI units.}
  \label{fig:pressure_pgv}
\end{figure*}

Frequency-dependent coupling was described by the H1 pressure-to-vertical-
velocity transfer-function estimator
\begin{equation}
 H_1(f)=\frac{S_{pv}(f)}{S_{pp}(f)},
 \label{eq:h1_transfer}
\end{equation}
where $S_{pv}(f)$ is the cross-spectral density between pressure $p$ and
vertical ground velocity $v$, and $S_{pp}(f)$ is the pressure autospectral
density. The H1 estimator treats pressure as input and vertical velocity as
output, so $|H_1|$ has units of m~s$^{-1}$~Pa$^{-1}$ and represents the
frequency-dependent seismic/acoustic coupling coefficient. Magnitude-squared
coherence was calculated alongside $H_1$. Continuous estimates used the
opening activity interval (Phase~I; see Section~\ref{sec:temporal_evolution});
event and continuous estimates overlap in time and were treated as
complementary rather than independent.

Event-window and continuous analyses both showed enhanced vertical response per
unit pressure near 4--5 and 20--24~Hz, with a weaker 13--14-Hz feature
(Figure~\ref{fig:transfer}). The H1 pressure-to-vertical-ground-velocity
transfer-function magnitude reached approximately
$5$--$7\times10^{-6}$~(m~s$^{-1}$)~Pa$^{-1}$ in the principal coupling bands,
with elevated coherence across parts of the same frequency ranges.

\begin{figure*}[!tbp]
  \centering
  \includegraphics[width=0.94\textwidth]{publication_pressure_dhz_transfer_and_coherence.pdf}
  \caption{Empirical frequency-dependent seismic/acoustic coupling coefficient
  at BCHH. H1 pressure-to-vertical-ground-velocity magnitude is shown for the
  high-quality event ensemble and continuous Phase~I, together with
  magnitude-squared coherence.}
  \label{fig:transfer}
\end{figure*}


The focused upper-stage analysis also identified a candidate seismic arrival
before the predicted ground-coupled airwave (Supplementary Section~S9).
Median persistent causal threshold crossings occurred at 1.863 and 1.895~s
in the 0.8--3- and 3--8-Hz bands, respectively, relative to the adopted
351~m~s$^{-1}$-reduced upper-stage time coordinate; the predicted acoustic
arrival at the seismometer was approximately 4.046~s. At 1.420~km range,
these onsets imply conditional path-average apparent speeds of approximately
0.75--0.76~km~s$^{-1}$. Their agreement across adjacent bands and the
amplitude-gated radial-plane polarization are suggestive of a direct
ground-propagating contribution. However, the source-time reference is
inferred, upstream response correction and preprocessing are not entirely
causal, and a single seismic station cannot uniquely identify P, S, or
Rayleigh waves. A seven-event stack did not corroborate a repeatable early
arrival (Supplementary Section~S9).

\section{Acoustic energetics}
\label{sec:energetics}

For the four named opening signals, acoustic exposure was integrated over
manually reviewed windows. For a progressive acoustic wave, intensity is
proportional to $p^2/(\rho c)$; assuming hemispherical spreading and no
atmospheric attenuation gives
\begin{equation}
 E_{2\pi}=\frac{2\pi r^2}{\rho c}\int p^2(t)\,dt,
 \label{eq:acoustic_energy}
\end{equation}
following standard acoustic-energy treatments
\citep[e.g.,][]{johnson_aster_2005}. The model does not explicitly correct
receiver ground reflection, source directivity, or local obstacles; changing
$2\pi$ to $4\pi$ does not resolve those effects. Acoustic TNT equivalence is
$E_{2\pi}/(4.184\times10^6~\mathrm{J\,kg^{-1}})$, an energy-unit conversion
rather than a blast-yield estimate. No confirmed regional Falcon~9 branch with
a stable dominant period was available, so efficiency-dependent yield scenarios
are confined to the Supplementary Material.

The continuous sequence-rate view in
Figure~\ref{fig:sequence_energy_rates} complements the event-window
measurements below by showing when acoustic energy and the local ground-motion
proxy accumulated, without treating every catalogue pick as an independent
energy release.

Within the processed infraBSU passband (approximately 0.05--80~Hz), the
manually reviewed principal-explosion window had a median positive-peak sound
pressure level of 156.8~dB re 20~$\mu$Pa and a median
hemispherical-equivalent radiated acoustic energy of
$2.2\times10^9$~J across DD1--DD3
(range $1.93$--$2.51\times10^9$~J). Expressed only as an energy-unit
conversion, this is equivalent to approximately 525~kg TNT.

In the separate de-duplicated sequence analysis, 51 selected event windows
merged into 47 nonoverlapping complexes. Their cumulative median radiated
acoustic energy was $2.27\times10^9$~J
(range $1.95$--$2.55\times10^9$~J), equivalent to 543~kg TNT
(range 467--610~kg). The cumulative value covers only the selected
measurement intervals rather than the complete 26-min sequence. A separate
source of uncertainty is the approximately $\pm20\%$ reproducibility of the
pressure calibration procedure; because acoustic energy scales with squared
pressure, this corresponds to a median-energy range of approximately
$(1.4\text{--}3.2)\times10^9$~J.

These acoustic-energy estimates can be compared with, but should not be
equated to, the chemical-energy reservoir of the launch vehicle. Published
Falcon~9 v1.2 specifications give nominal propellant masses of approximately
411~t for the first stage and 111.5~t for the second stage, for approximately
522.5~t of LOX/RP-1 at full load
\citep{<Falcon9_datasheet_key>}. Using the nominal LOX/RP-1 mixture ratio and
the heating value of RP-1 gives a full-load chemical-energy reservoir of order
$6$--$7\times10^{12}$~J. The exact propellant inventory aboard AMOS-6 at the
time of failure is not publicly documented, and the accident occurred during
the propellant-loading sequence; the full-load value is therefore an upper
bound rather than an estimate of the instantaneous onboard energy.


\section{Regional detectability and atmospheric propagation}
\label{sec:regional}

Regional screening used broad physically plausible acoustic-celerity and
seismic-velocity windows. Pressure waveforms were band-pass filtered from 0.2
to 1.0~Hz and evaluated within apparent-celerity windows of
180--350~m~s$^{-1}$, subdivided for candidate generation into nominal
thermospheric-or-slow (180--260~m~s$^{-1}$), stratospheric
(260--310~m~s$^{-1}$), and tropospheric (310--350~m~s$^{-1}$) windows. These
labels do not identify actual atmospheric propagation paths. For each pressure
phase window, the maximum absolute pressure was compared with equal-duration
pre- and post-event controls. Robust background noise was
$\sigma=1.4826\,\mathrm{MAD}$. Candidates required SNR $\geq2.5$, an event
peak at least 1.25 times the larger control peak, complete waveform coverage,
and a 120-s guard between controls and the expected-arrival interval. Candidate
arrival time was the maximum Hilbert-envelope amplitude within the phase
window. These permissive criteria generated candidates for visual and
network-coherence review rather than formal detections.

Regional screening began with 171 station--sensor targets, of which 165
produced usable waveform records. The current run flagged 94 candidate phase
windows at 68 stations. An earlier candidate set was visually reviewed; the
current exported review table does not retain those decisions, so completion
of manual classification for all current windows is not asserted here.

Three prominent 0.2--1.0-Hz pressure packets at N4.257A.EP.BDF,
N4.S54A.EP.BDF, and N4.R55A.EP.BDF (381--1079~km) stand above the surrounding
background and are retained as possible regional arrivals from the Falcon~9
failure. Their peak amplitudes are 0.162--0.352~Pa and apparent celerities
192.7--224.7~m~s$^{-1}$. The packets at S54A and R55A are particularly
distinctive; 257A contains more competing transient activity. Their occurrence
within the slow-celerity screening window does not establish a thermospheric
path. Source association remains uncertain because coherent network moveout
has not been established. They are not used for period--yield analysis.
Waveform panels, maps, and screening details are in Supplementary Section~S8.

The upper-air observations showed strongly height-dependent winds, with
southerly to southwesterly flow through much of the lower troposphere, stronger
westerly flow near 10--15~km altitude, and a reversal to easterly flow in the
lower stratosphere. The AMPS temperature profile showed decreasing sound speed
through the troposphere followed by recovery above the tropopause, producing a
strongly height- and azimuth-dependent effective sound-speed structure. These
conditions are consistent with substantial directional refraction and possible
regional acoustic shadowing, but no atmospheric propagation model was used to
predict specific return ranges.


\section{Discussion}
\label{sec:discussion}

\subsection{Forensic chronology and the precursor question}

The BCHH record supplies a chronology independent of vehicle telemetry. One
question posed during the initial forensic review was whether the record
contained an external precursory acoustic or seismic signal that might indicate
an externally initiated event. The 48-h six-channel review found no such signal
before the initial upper-stage explosion, consistent with the November~2016
forensic report \citep{thompson_mcnutt_fireball_2016} and the 2017 AGU
presentation \citep{thompson_agu_falcon9_2017}. This result does not exclude
weak, internal, non-impulsive, or poorly coupled processes, but it provides no
geophysical evidence for a conspicuous external precursor.

The opening chronology is now constrained without assigning absolute UTC to the
video. The first event is common to the BCHH record and the public
USLaunchReport recording and defines the reference. Later source intervals are
measured within that same public recording, while later arrivals are measured
independently on DD2 and DD3. Proprietary SpaceX views, including UCS-3 near
BCHH and cameras at SLC-40, provide higher-frame-rate validation without being
mixed into the public-video clock. This is a stronger construction than
comparing BCHH arrivals with source times obtained by reducing those same
arrivals at 351~m~s$^{-1}$. SpaceX later attributed the
initiating failure to a composite overwrapped pressure vessel in the
second-stage liquid-oxygen tank \citep{spacex_anomaly_2017,nasa_sma_falcon9}.
BCHH cannot identify that internal mechanism. Its abrupt onset is compatible
with, but does not prove, the reported rapid transition from anomalous telemetry
to second-stage loss.

\subsection{Differential timing and weak-shock propagation}

The independent event intervals provide the strongest evidence for
propagation faster than the nominal acoustic speed. The public recording gives
a 0.07--0.14-s differential reduction for the principal explosion, with
substantially smaller reductions for the payload events; the proprietary UCS-3
record independently gives the same ordering and scale. The principal range
reflects the preferred 0--2-frame visual-onset uncertainty. Agreement between
independent camera views is more important here than small differences between
their frame-scale residuals.

This comparison does not require the camera clock to be synchronized to BCHH.
Nor, for events at the same nominal source position, does it depend strongly on
the exact value of the 351~m~s$^{-1}$ reference speed: the common linear travel
time cancels when event intervals are differenced. The reference speed remains
useful for catalogue reduction and for expressing nominal predicted arrivals.
For the payload events, source position is less certain, and a displacement of
only several metres begins to matter for residuals of a few hundredths of a
second. Their timing is therefore supporting rather than decisive evidence.

The amplitude dependence strengthens this interpretation: the weak-shock
sensitivity calculation reproduces the observed ordering and approximate scale
without using the video timing to set its parameters. The agreement should not
be overinterpreted, because extrapolating $1/r$ pressure growth inward to 10~m
is a sensitivity calculation rather than a source-pressure inversion.

The integrated source-to-receiver timing should also not be confused with the
local array-crossing velocity. The preferred principal-event array solution is
352.5~m~s$^{-1}$, close to the meteorological effective sound speed at BCHH.
The measured positive overpressure corresponds to only $M\approx1.006$ if
interpreted as a normal-shock jump. These local observations are consistent
with a wave that was only weakly nonlinear by 1.4~km. They do not argue against
faster propagation earlier on the path: a blast wave decelerates as its
overpressure decreases. Conversely, the approximately 369~m~s$^{-1}$
short-window Bartlett value would require an overpressure of order 13~kPa if
interpreted literally as the local shock speed, and is therefore unlikely to
represent the true wavefront velocity at BCHH.

Waveform morphology is compatible with this blast-wave interpretation but
does not independently establish Mach number. Taken together, the differential
timing, measured pressure, local array speed, waveform shape, and weak-shock
sensitivity calculation favor a nonlinear blast that propagated supersonically
closer to the vehicle and decayed toward ordinary acoustic propagation before
reaching BCHH. The data do not require a strongly supersonic front at the
receiver.

\subsection{Seismic/acoustic coupling and energy partition}

One of the more transferable results is the empirical local
seismic/acoustic coupling coefficient. Pressure scaled nearly linearly with
three-component vector PGV across the measurement-quality subset, and removing
the principal measurement group changed the fitted relation only modestly. The
result therefore does not arise solely from the largest event, although it
remains an empirical BCHH relation rather than a universal
pressure-to-ground-motion conversion. This velocity-to-pressure ratio is
directly comparable in form to seismic/acoustic coupling coefficients derived
from Hunga Tonga--Hunga Ha'apai pressure waves
\citep{anthony_hunga_coupling_2023,macpherson_hunga_coupling_2025}, while
remaining conceptually distinct from energy-partition metrics such as the
volcano acoustic-to-seismic ratio (VASR) \citep{johnson_aster_2005}.

The H1 analysis reveals frequency dependence hidden by the broadband
peak-amplitude relation, with repeatable enhanced response in the bands
identified above. Its spectral magnitude need not equal the broadband coupling
coefficient because H1 uses vertical motion and spectral averages within
individual frequency bands. Agreement between event-window and continuous
estimates supports a repeatable local seismic/acoustic response, comparable in
scale to controlled-explosion, volcanic, and Hunga Tonga--Hunga Ha'apai
measurements \citep{novoselov_ground_coupling_2020,anthony_hunga_coupling_2023,
macpherson_hunga_coupling_2025}. A second seismic station and a shallow
velocity model would be needed to relate this response more directly to local
material properties and subsurface structure.

The pressure calibration remains an important limitation on both the
coupling and energetic estimates, but the lift-test calibration is supported
by the independent launch comparison. For the likely 1-inch DD3 head, its
nominal sensitivity of 9.2~counts~Pa$^{-1}$ is close to the adopted
10.0~counts~Pa$^{-1}$ lift-test value: using the nominal value would raise
DD3-derived pressures by only 8.7\%. This agreement does not establish an
absolute laboratory calibration, and the anomalously low DD2 sensitivity
remains an important qualification, but it supports use of the empirical
factors within the stated $\pm20\%$ pressure uncertainty.

The principal explosion accounts for nearly all of the selected
acoustic-energy budget under the adopted processing and $2\pi$ radiation
model. These quantities are finite-passband, model-equivalent radiated acoustic
energies rather than blast yields. Absolute calibration, possible sensor
over-range response, baseline choice, receiver reflection, source directivity,
and local obstacles are correspondingly more important than the reported
numerical precision.

A useful interpretation requires separating three energy scales:
the chemical energy stored in the propellants, the fraction participating
promptly in the pressure-producing blast, and the still smaller fraction
radiated acoustically. Nominal Falcon~9 Full Thrust propellant capacity
corresponds to a full-load chemical-energy reservoir of order
$6$--$7\times10^{12}$~J. The measured principal-explosion acoustic energy is
therefore only about $3\times10^{-4}$, or 0.03\%, of that reservoir. This
chemical-to-acoustic ratio spans both conversion steps and should not be
interpreted as an acoustic efficiency. Moreover, the vehicle failed during
propellant loading, so the actual onboard chemical-energy reservoir was below,
or at most equal to, the nominal full-load value.

The distinction is illustrated by the 2026 New Glenn static-fire explosion.
Regional infrasound observations gave an effective blast yield of approximately
0.131~kt TNT, or $5.5\times10^{11}$~J, and comparison with plausible onboard
chemical-energy reserves implied a blast-coupling fraction of approximately
3--4\% for a full-fill scenario
\citep{silber_scamfer_newglenn_2026}. This is the fraction of chemical energy
represented by the inferred impulsive blast, not the fraction radiated as
acoustic energy. The New Glenn study therefore constrains a different step in
the energy partition from the close-range Falcon~9 measurement.

As a sensitivity comparison, applying the same 3--4\% blast-coupling fraction
to a fully loaded Falcon~9 would give an effective impulsive blast energy of
approximately $(2.0$--$2.7)\times10^{11}$~J, or about 48--64~t TNT
equivalent. The measured $2.2\times10^9$-J radiated acoustic energy would then
correspond to approximately 0.8--1.1\% of that effective blast energy. These
values are not an independent Falcon~9 yield estimate: neither the actual
propellant fill fraction nor the chemical-to-blast coupling is known. They do,
however, demonstrate that the apparently small 0.03\% chemical-to-acoustic
ratio does not require an exceptionally inefficient acoustic source. A
two-stage partition,
\[
 E_{\mathrm{chemical}}
 \rightarrow E_{\mathrm{blast}}
 \rightarrow E_{\mathrm{acoustic}},
\]
with a few percent entering the impulsive blast and of order one percent of
that blast energy appearing in the present radiated-acoustic estimate, is
consistent with both the Falcon~9 measurements and the New Glenn comparison.
The detailed sensitivity calculation is given in Supplementary Section~S7.

\subsection{Regional propagation and candidate arrivals}

No regional counterpart to the Falcon~9 explosion was confirmed. The
permissive screen produced numerous candidates, but manual review and the
northern record section found no coherent regional moveout; the three
prominent pressure packets remained possible arrivals with uncertain source association. The event therefore
cannot support a regional period--yield analysis of the kind applied to New
Glenn \citep{silber_scamfer_newglenn_2026}.

The absence of a confirmed regional association is not necessarily inconsistent with the strong
close-range signal. Contemporaneous KSC wind-profiler and AMPS sounding data
show a strongly height- and azimuth-dependent atmosphere, including decreasing
temperature-dependent sound speed through the troposphere, strong westerly
flow near 10--15~km, and a wind reversal in the lower stratosphere. These
observations indicate conditions favorable for substantial directional
refraction and are consistent with possible acoustic shadowing or strongly
favored propagation corridors. Because the measured sounding extends only
through the lower stratosphere and no ray-tracing or full-wave propagation
calculation is included here, we do not use the atmospheric structure to
predict specific return ranges or explain individual station
non-detections. Wind-dependent spectral filtering provides an additional
reason not to transfer regional period--yield relations without propagation
modelling \citep{silber_wind_filtering_2026}.

\subsection{Monitoring implications and limitations}

A single stand-off site constrained onset, chronology, relative amplitude,
source direction, acoustic speed, and local seismic/acoustic coupling while pad
access was restricted. A permanent KSC seismo-acoustic network could build on
that result. Several compact pressure arrays paired with separated
three-component seismic stations would provide triangulation and path sampling,
separate direct ground waves from local seismic/acoustic conversion, and test
whether the BCHH seismic/acoustic coupling coefficient transfers between sites. Continuous
recording would preserve the background needed for both real-time situational
awareness and later forensic analysis.

The main limitations are the single site, metre-scale reconstructed sensor
positions, empirical pressure calibration, pressures beyond nominal sensor
ranges, manual event identification, uncertain absolute video timing, and
first-order radiation assumptions. The provisional $\pm20\%$ pressure
calibration uncertainty affects all absolute pressure-dependent results.
Baseline, geometry, propagation, and possible over-range nonlinearity add
uncertainties outside that percentage. Nearby-launch raw-count checks support
relative DD2/DD3 stability, but do not identify DD2's nominal head or establish
an absolute calibration. The timing, direction, and scaling results remain
useful; absolute energy, yield, and seismic-phase interpretations require more
caution.


\section{Conclusions}
\label{sec:conclusions}

The BCHH array captured an unusually complete close-range record of the 2016
Falcon~9 AMOS-6 failure. A gap-free six-channel review found no precursor
signal during the 37~h before the initial upper-stage event, although weak,
internal, non-impulsive, or poorly coupled processes remain outside that
constraint. The subsequent 26-min sequence contains 153 accepted multichannel
pressure transients distributed among four prominent activity phases.

Independent public-video and BCHH timing shows that the principal explosion
gained about 0.07--0.14~s of propagation time relative to the initial event
over the 1.4-km path, while the smaller payload events show smaller positive
residuals. Higher-frame-rate proprietary SpaceX recordings independently
support the event identifications and the same pressure-dependent ordering,
which is also reproduced in scale by the weak-shock sensitivity calculation.

The principal event did not remain strongly supersonic to BCHH. Its preferred
local array-crossing speed was 352.5~m~s$^{-1}$ and its 1.39-kPa positive
overpressure corresponds to only $M\approx1.006$ under a normal-shock
interpretation. The combination is consistent with a stronger nonlinear blast
closer to the vehicle that decayed toward the ambient acoustic speed by
1.4~km. Its strongest pressure waveform is Friedlander-like, but waveform
shape alone is not used to establish the Mach number.

Within the processed infraBSU passband, the principal explosion reached a
median positive-peak sound pressure level of 156.8~dB re 20~$\mu$Pa and a
median hemispherical-equivalent radiated acoustic energy of
$2.2\times10^9$~J. That energy is distinct from explosive yield. A fully loaded
Falcon~9 contains of order $6$--$7\times10^{12}$~J of propellant chemical
energy, so the measured acoustic energy is only about 0.03\% of that nominal
reservoir and spans both chemical-to-blast and blast-to-acoustic conversion.

Array back azimuths independently support SLC-40 as the source. Event
pressures scaled nearly proportionally with vector PGV, giving a median
seismic/acoustic coupling coefficient of
$2.3\times10^{-6}$~(m~s$^{-1}$)~Pa$^{-1}$, while H1 analysis identified
repeatable coupling bands near 4--5 and 20--24~Hz. No coherent regional
counterpart was confirmed despite systematic screening and manual review.
Contemporaneous upper-air observations indicate strongly direction-dependent
propagation, but detailed regional ray paths were not modelled.

These observations show the value of close-range, continuously recording
seismo-acoustic instrumentation for reconstructing complex launch-vehicle
failures when engineering telemetry and direct site access are limited. A
permanent distributed KSC seismo-acoustic network would improve source
triangulation, constrain blast-wave decay, and provide repeat measurements of
local seismic/acoustic coupling. The archived waveforms, direct arrival picks,
same-camera frame-timing metadata, catalogue, and reproducible workflow provide
a reference dataset for that monitoring design and for future studies of
complex atmospheric explosions.

\begin{acknowledgements}
We thank Michael Wagner and USLaunchReport for recording the accident and for
written permission to use selected material in presentations and publications.
We thank SpaceX for providing additional camera recordings used to check the
opening chronology.
We thank Kennedy Space Center personnel for facilitating the BCHH deployment
and access to meteorological observations. Regional waveform data were
accessed through FDSN services; the IU Global Seismographic Network is cited
in the reference list \citep{asl_gsn_iu_1988}.
\end{acknowledgements}


\section*{Funding}

The present reanalysis and manuscript received no dedicated funding. SpaceX
provided US\$5,000 to USF in September 2016 for the original forensic analysis
and report. A contemporaneous US\$25,000 Florida Space Grant Consortium award
supported the broader research program, comprising US\$12,500 from the
Consortium and US\$12,500 in USF matching funds. Neither funder influenced the
present reanalysis, interpretation, manuscript, or decision to submit.


\section*{Data and code availability}

The Python notebooks, project modules, shared configuration, environment
specification, data manifest, output conventions, and derived products will be
archived as a versioned release in a repository issuing a persistent DOI. BCHH
waveforms, metadata, calibration information, manual picks, meteorological
inputs, event catalogue, and derived tables will be deposited in the same or a
linked record. The 48-h six-channel input used for the precursor screen is
separate from the short accident-analysis record and will be identified with
its exact coverage and processing provenance. The DOI and repository URL will be inserted before submission.

Regional waveforms remain available from their FDSN data centres and network
operators. The archive will include the station/channel manifest, availability
audit, screening metrics, and manual-review status/decisions. Selected regional
waveforms used for candidate validation will be retained as MiniSEED with
event-epoch StationXML and exact request/coverage metadata. The November~2016
technical report is cited as an unpublished historical source and is not
publicly available \citep{thompson_mcnutt_fireball_2016}; the 2017 AGU
presentation abstract is publicly archived \citep{thompson_agu_falcon9_2017}. The USLaunchReport and
Scott Manley videos remain at their cited URLs
\citep{uslaunchreport_amos6_video_2016,manley_spacex_explosion_2016}. Copyright
in the source footage remains with USLaunchReport.com/Michael Wagner; the
archive will provide synchronization metadata and code rather than redistribute
the video. The proprietary SpaceX camera files cannot be redistributed; the archive
will instead include the reviewed UCS-3/SLC-40 frame and timecode observations,
together with code used to calculate intervals within each camera separately.
The public USLaunchReport PTS picks used for the primary timing analysis will
also be archived.


\section*{Competing interests}

The authors declare no competing interests. SpaceX support was limited to the
original 2016 forensic work; SpaceX had no role in the present reanalysis,
interpretation, manuscript preparation, or submission decision.

% The Seismica class selects abbrvnat_seismica_upcasetitle automatically.
\bibliography{references}

\end{document}
